
\documentclass[11pt]{article}
\usepackage{amsmath,amssymb}
\usepackage[a4paper,margin=1in]{geometry}

\title{Introduction to the Cahn--Hilliard Equation}
\date{}

\begin{document}
\maketitle

\section{Physical Idea}

The Cahn--Hilliard equation describes phase separation in a binary mixture while conserving mass.
Let
\[
c(\mathbf{x},t)
\]
denote the concentration of one component.

Mass conservation requires
\[
\frac{d}{dt}\int_\Omega c\,dV = 0.
\]

\section{Free Energy Functional}

The free energy is
\[
F[c]
=
\int_\Omega
\left(
f(c)
+
\frac{\kappa}{2}|\nabla c|^2
\right)
dV.
\]

The bulk free energy is often chosen as
\[
f(c)=\frac{1}{4}(c^2-1)^2,
\]
which has minima at \(c=\pm1\).

\section{Chemical Potential}

The chemical potential is the variational derivative of the free energy:
\[
\mu=\frac{\delta F}{\delta c}.
\]

This yields
\[
\mu=f'(c)-\kappa\nabla^2 c.
\]

For the double-well potential,
\[
f'(c)=c^3-c,
\]
so
\[
\mu=c^3-c-\kappa\nabla^2 c.
\]

\section{Mass Conservation and Flux}

The conservation law is
\[
\frac{\partial c}{\partial t}+\nabla\cdot\mathbf{J}=0.
\]

Assuming a diffusive flux
\[
\mathbf{J}=-M\nabla\mu,
\]
where \(M\) is the mobility, gives
\[
\frac{\partial c}{\partial t}
=
\nabla\cdot(M\nabla\mu).
\]

\section{Classical Cahn--Hilliard Equation}

Combining the above relations gives
\[
\frac{\partial c}{\partial t}
=
\nabla\cdot
\left(
M\nabla
\left(
f'(c)-\kappa\nabla^2 c
\right)
\right).
\]

For constant mobility,
\[
\frac{\partial c}{\partial t}
=
M\nabla^2
\left(
f'(c)-\kappa\nabla^2 c
\right).
\]

Using the quartic potential,
\[
\frac{\partial c}{\partial t}
=
M\nabla^2(c^3-c-\kappa\nabla^2 c).
\]

\section{Fourth-Order Nature}

Since
\[
\nabla^2(\nabla^2 c)=\nabla^4 c,
\]
the equation is fourth order in space.

A common finite-element formulation introduces
\[
\mu=f'(c)-\kappa\nabla^2 c,
\]
and solves
\[
\frac{\partial c}{\partial t}
=
\nabla\cdot(M\nabla\mu).
\]

\section{Energy Dissipation}

The free energy decreases monotonically:
\[
\frac{dF}{dt}
=
-\int_\Omega M|\nabla\mu|^2\,dV
\le 0.
\]

\section{Spinodal Decomposition}

For a nearly homogeneous state
\[
c=c_0+\epsilon,
\]
perturbations grow when
\[
f''(c_0)<0.
\]

This leads to spontaneous phase separation.

\section{Linear Stability Analysis}

The growth rate of a Fourier mode is
\[
\lambda(k)
=
-Mk^2
\left(
f''(c_0)+\kappa k^2
\right).
\]

Long wavelengths can grow when \(f''(c_0)<0\), while short wavelengths are stabilized by the gradient term.

\section{Boundary Conditions}

Typical no-flux conditions are
\[
\nabla\mu\cdot n=0,
\qquad
\nabla c\cdot n=0.
\]

\section{Relation to Allen--Cahn}

Allen--Cahn:
\[
\partial_t c=-M\mu.
\]

Cahn--Hilliard:
\[
\partial_t c=\nabla\cdot(M\nabla\mu).
\]

Allen--Cahn does not conserve mass, while Cahn--Hilliard does.

\section{Coupling to Mechanics}

A chemo-mechanical free energy can be written as
\[
F=
\int
\left(
f(c)
+\frac{\kappa}{2}|\nabla c|^2
+\psi_{\mathrm{elastic}}(\varepsilon,c)
\right)dV.
\]

The chemical potential becomes
\[
\mu
=
\frac{\partial f}{\partial c}
-\kappa\nabla^2 c
+
\frac{\partial \psi_{\mathrm{elastic}}}{\partial c}.
\]

Applications include alloys, battery electrodes, hydrides, and phase-field models.

\end{document}
